\scanpage{10}

\mspage{5)}
On utilise le weak lifting lemma pour

\margin{\uline{NB} On utilise $y\in C_{\rm Cof}$ \deleted{\unsure{en}} ici}

\diagram{\begin{tikzcd}[ampersand replacement=\&,column sep=small,row sep=small]
\varnothing \arrow[r] \arrow[d,tail] \& x \arrow[r,tail,dashed,"i","\sim"'] \arrow[d,"f","\sim"'] \& \bar x \arrow[dl,dashed,"{\deleted{\unsure{g}}\alpha}","\sim"'] \\
y \arrow[r,"1"] \& y \arrow[ur,bend right,dashed,"g"'] \& {}
\end{tikzcd}}

donc il existe $i,\deleted{\unsure{g}}\alpha$ avec $\deleted{\unsure{g}}\alpha i=f$, \deleted{\unsure{et}} $i\in W\cap\mathrm{cof}$, et $\alpha$ admet une section $g$ \inserted{$\alpha g=1_y$}.

Je dis que $(g,i)$ est l'inverse \inserted{cherché.} \deleted{\unsure{de $(f,\mathrm{id}_y)$}} \deleted{Considérons alors}

\diagram{\begin{tikzcd}[ampersand replacement=\&,column sep=small,row sep=small]
x \arrow[dd,"{1_x}"'] \arrow[dr,"{gf}"] \arrow[dddr,dashed,"{\deleted{\unsure{gf}}f}"'] \& x \arrow[d,tail,"i","\sim"'] \arrow[ddl,tail,"{1_x}","\sim"'] \\
{} \& \bar x \arrow[dd,dashed,"{\alpha}","\sim"'] \\
x \arrow[dr,dashed,"f"'] \& {} \\
{} \& y
\end{tikzcd}}

Pour $(gf,i)\sim\underline{1}_x$, on a envie d'utiliser le diagramme commutatif ci-contre, le seul ennui c'est que, on a $\alpha\in W$ mais pas néc. $\alpha\in\mathrm{cof}$, donc $\alpha i$ $(=f)$ pas néc. dans $\deleted{\unsure{W|\mathrm{cof}}}\mathrm{cof}$ (alors qu'il est dans $W$). Ça marcherait si on pouvait trouver $y\xrightarrow[\sim]{\beta}y'$ \inserted{dans $W$,} tel que le composé $\beta f$ soit non seulement dans $W$, mais de plus, dans cof.

Visiblement, ce n'est pas tjrs le cas (prendre le cas où les cof. sont les monos, et où $f$ n'est \uline{pas} un mono !). Ça se présente mal !

\diagram{\begin{tikzcd}[ampersand replacement=\&,column sep=small,row sep=small]
x \arrow[d,"f"'] \arrow[ddr,tail] \& {} \\
y \arrow[dr,"\beta"'] \& {} \\
{} \& y'
\end{tikzcd}}

On va alors plutôt définir d'abord le foncteur canonique

\[
\begin{aligned}
 \widetilde C&\longrightarrow W^{-1}C\\
 &\unsure{\varphi|\colon|{}}\ x\longmapsto x\\
 &(f,i)\longmapsto [i]^{-1}[f]
\end{aligned}
\]

il est immédiat que c'est compatible avec les identités et la composition. Donc on a bien foncteur, \deleted{surjectif} \inserted{bijectif} sur objets. Je dis que c'est surjectif sur les flèches ; Il suffit de voir, par composition, que si $f:x\longrightarrow y$ est dans $W$, alors $[f]^{-1}$ provient de $\widetilde C$. On va prendre \deleted{$(y,x)$} $y\longrightarrow x$ de $\widetilde C$, défini par $(g,i)$, \deleted{À prouver que} \inserted{i.e.} $[i]^{-1}[g]$, pour prouver que $[i]^{-1}[g]=[f]^{-1}$, il suffit de prouver que $[f][i]^{-1}[g)=[1_y]$, or $\alpha i=f$ entraîne $[\alpha][i]=[f]$ d'où $[f][i]^{-1}=[\alpha]$ donc le pr. mb est $[\alpha][g]=[\alpha g]=[1_y]$, OK.
% archive sheet number: 9
